\documentclass[]{article}

\usepackage{amsmath}
\usepackage{multirow}
\usepackage{natbib}
\usepackage{graphicx} 


\begin{document}

Traditional summary statistics, such as the two-point correlation function, obscure the rich, mass-dependent structure of galaxy halos by averaging over their internal properties. We present a framework that bypasses this information loss by directly modeling the three-dimensional positions of galaxies as a mass-conditioned spatial point process. Applying a Neyman-Scott process model to a suite of ten synthetic galaxy catalogs, we perform a maximum likelihood estimation to recover the underlying Halo Occupation Distribution (HOD) parameters that govern satellite populations. Our model recovers the input HOD parameters with a small, well-understood systematic bias. Using the Akaike Information Criterion for model selection, we find decisive evidence that the satellite radial concentration increases with host halo mass, revealing a subtle break in the self-similarity of halo structure. Furthermore, by employing a marked correlation function with luminosity as the mark, we quantify the spatial segregation within halos, finding that more luminous galaxies are preferentially located near halo centers. A residual analysis precisely quantifies the breakdown of the 1-halo model at scales of 5-10 Mpc/h, where inter-halo clustering becomes the dominant contribution. This work demonstrates that direct likelihood-based modeling of spatial point patterns can extract detailed astrophysical information from galaxy catalogs, providing a powerful alternative to traditional summary statistics for analyzing next-generation cosmological surveys.

\end{document}
                